%% title: [A] Abstention under shift is a composition, not a score
%% target: Sections/Results.tex
%% anchor: The representation is the expensive and robust component; the classifier is the cheap one where interventions belong, and a new camera needs a new classifier, which takes minutes and no labels, rather than a new model.
%% action: after
%% why: Backing off is what makes a model usable under shift: thresholds fitted on target images keep 95 % precision, self-training halves abstention and adds 10 points of correct answers, poor images are handled without being told, and novelty remains the weak link.
%% source: /data/au761367/paper_numbers/22.md (dev/097 on saved campaign predictions; traps: 15,141 known-label probe crops, 11 halvings; test fold: 632,913 images, 11 halvings)
%% check: Rewritten 2026-11-09 with campaign models. Cut: expected utility (loses for every head but I in the rerun), the genus-target and temperature paragraphs (margin heads only), B8 / P5 / B3rep5x rows (no rerun). | Paste after step 20; it closes the results (step 21 is dropped).
%% status: ready
\subsection{Backing off under domain shift}\label{sec:results_backoff}
On the test fold the greedy policy of \cref{sec:backoff} answers 99.2 to 99.7\,\% of images at 95\,\% precision, and the genus and family answers carry their target because each threshold was fitted on the images that reach its rank.
On held-out trap nights the same policy abstains on 17 to 21\,\% of crops for models trained without trap images, and self-training cuts this to about 7\,\% and raises the share answered correctly by about 10 points, mostly through species answers (\cref{tab:backoff}).

\begin{table}[htb]
    \centering
    \caption{Share of images answered correctly at the rank returned, greedy policy at 95\,\% precision per rank. Test fold: all images, eleven random halvings. Trap: probe crops with a known label, thresholds fitted on half of the (trap, night) groups, mean of eleven halvings. One seed per model.}
    \label{tab:backoff}
    \begin{tabular}{lcc}
        \toprule
        model & test fold & trap \\
        \midrule
        EfficientNetV2-S, F & -- & 75.5\,\% \\
        EfficientNetV2-S, BU & 94.3\,\% & 77.2\,\% \\
        EfficientNetV2-S, I & -- & 76.8\,\% \\
        \quad BU + self-training & 94.3\,\% & \textbf{88.6}\,\% \\
        DINOv3 ConvNeXt-L, BU & \textbf{94.7}\,\% & 79.0\,\% \\
        \quad BU + self-training & -- & 88.0\,\% \\
        \bottomrule
    \end{tabular}
\end{table}

\emph{Poor images.} From the sharpest to the blurriest fifth of trap crops, the BU model's species answers fall from 61 to 22\,\% and its abstentions rise from 4 to 42\,\%, while species precision stays at 0.891 on the blurriest fifth (DINOv3: 75 to 25\,\%, 4 to 40\,\%, 0.924).
The policy backs off on poor images without being told about image quality.

\emph{Thresholds must be fitted on target images.} Fitted on the test fold and applied to trap images, the policy misses its 95\,\% target by 7 to 17 points.

\emph{Unseen species.} No back-off rule fixes novelty.
Of the trap crops of unseen species in known genera, 35\,\% receive a species answer, all wrong; on museum-style images the policy names a species for 69\,\% of unseen-species images.
Detecting novelty (\cref{sec:results_openset}), not choosing the rank, is the weak link under shift.
